%======================================================================
% CAPÍTULO 04: PRESTACIONES SOCIALES
%======================================================================
% Anticipos, finiquitos, judiciales, movimientos y órdenes de pago
% Fundamento Legal: Art. 57, 60, 61 de la Ley SSSIFANB
%======================================================================

\chapter{Módulo de Prestaciones Sociales}
\label{ch:prestaciones}

\section{Descripción General}

El Módulo de Prestaciones Sociales es el ``corazón fiduciario'' del sistema, responsable operativo del Capítulo IV de la \leynegroprimero. Sus algoritmos matemáticos resguardan miles de millones de bolívares al año y dictaminan el patrimonio de una vida de servicio basándose directamente en el tiempo de cotización validado desde el \hyperref[ch:afiliacion]{Módulo de Afiliación}.

\section{Acceso al Módulo}

Para acceder al módulo de Prestaciones:

\begin{enumerate}[leftmargin=*]
    \item Seleccione \textbf{Prestaciones} en el menú lateral
    \item Se mostrará el menú de opciones disponibles
\end{enumerate}

\begin{figure}[H]
\centering
\begin{minipage}{0.48\textwidth}
\centering
\includegraphics[width=\textwidth]{5-prestaciones-menu.png}
\caption{Módulo de Prestaciones}
\label{fig:prestaciones-menu}
\end{minipage}
\hfill
\begin{minipage}{0.48\textwidth}
\centering
\end{minipage}
\end{figure}

\section{Submódulos Disponibles}

El módulo de Prestaciones Sociales se organiza en los siguientes submódulos funcionales:

\begin{itemize}[label=\textcolor{sandraGreen}{\faIcon{angle-right}}, leftmargin=*]
    \item \textbf{\faIcon{money-bill-wave} Anticipos}: Solicitud de percepción anticipada de fondos fiduciarios.
    \item \textbf{\faIcon{file-signature} Finiquitos}: Liquidación final integral de prestaciones al cese del servicio.
    \item \textbf{\faIcon{balance-scale} Judiciales}: Gestión de medidas cautelares y embargos dictaminados por tribunales.
    \item \textbf{\faIcon{history} Movimientos}: Historial contable de operaciones del fideicomiso.
    \item \textbf{\faIcon{file-invoice-dollar} Órdenes de Pago}: Control estricto de desembolsos a cuentas del militar.
    \item \textbf{\faIcon{layer-group} Lotes}: Procesamiento asíncrono masivo para actualización de expedientes.
    \item \textbf{\faIcon{chart-bar} Reportes}: Generación de proyecciones y estados de cuenta certificados.
\end{itemize}

\section{Fundamento Matemático - Asignación de Antigüedad (Art. 57)}

\keyconcept{Motor Lógico del Fideicomiso y Prestaciones}\\
El ciclo asincrónico procesa contablemente los pagos por antigüedad e invoca las siguientes reglas de bloque comercial fiduciario de la ley:

\begin{lstlisting}[language=Java, caption=Cálculos Fiduciarios de Antigüedad]
// 1. Corridas Trimestrales de Garantía
function procesarGarantiaTrimestral(sueldoIntegralDate) {
    const sueldoDiario = sueldoIntegralDate / 30;
    return sueldoDiario * 15; // Abonar a Fideicomiso
}

// 2. Corridas Anuales (Triggers después de Años_Servicio > 1)
function procesarBonoAdicional(sueldoIntegralDate, anosServicio) {
    if (anosServicio > 1) {
        const sueldoDiario = sueldoIntegralDate / 30;
        return sueldoDiario * 2; // Extra por cada año vencido
    }
    return 0;
}
\end{lstlisting}

\section{Anticipos del Fideicomiso}

Los anticipos permiten al afiliado disponer de una porción de sus prestaciones acumuladas antes de la desvinculación definitiva.

\subsection*{Tipos de Anticipo}

\begin{longtable}{p{3cm}p{4cm}p{5cm}}
\toprule
\textbf{Tipo} & \textbf{Monto Máximo} & \textbf{Restricciones} \\
\midrule
\endhead
\bottomrule
\caption{Tipos de anticipo disponibles}
\label{tab:tipos-anticipo}
\endlastfoot

Educativo & 30\% & Para estudios del afiliado o dependientes \\
Médico & 40\% & Para tratamientos de salud \\
Vivienda & 50\% & Para compra o mejora de vivienda \\
Emergencia & 25\% & Para situaciones imprevistas calificadas \\
\end{longtable}

\subsection*{Requisitos para Solicitar}

\begin{enumerate}[leftmargin=*]
    \item Tener al menos 2 años de cotización continua
    \item No tener anticipos pendientes de liquidación
    \item Cumplir con los requisitos específicos del tipo de anticipo
    \item Tener cuenta bancaria activa registrada en el sistema
\end{enumerate}

\begin{tcolorbox}[
    enhanced,
    colback=sandraSoft,
    colframe=sandraDark,
    width=0.9\textwidth,
    halign=center
]
\small
\textbf{Capacidad de Pago}\\
\smallskip
El sistema verifica que la suma de anticipos solicitados no supere el 70\% del total de las prestaciones acumuladas. Este cálculo protege la integridad de los fondos de fideicomiso.
\end{tcolorbox}

\section{Lista y Solicitud de Anticipos}

\begin{figure}[H]
\centering
\begin{minipage}{0.48\textwidth}
\centering
\includegraphics[width=\textwidth]{5.1-lista-de-anticipos.png}
\caption{Lista de anticipos}
\label{fig:lista-anticipos}
\end{minipage}
\hfill
\begin{minipage}{0.48\textwidth}
\centering
\includegraphics[width=\textwidth]{5.2-nuevo-anticipo.png}
\caption{Nuevo anticipo}
\label{fig:nuevo-anticipo}
\end{minipage}
\end{figure}

\begin{figure}[H]
\centering
\begin{minipage}{0.48\textwidth}
\centering
\includegraphics[width=\textwidth]{5.3-solicitar-anticipo.png}
\caption{Solicitar anticipo}
\label{fig:solicitar-anticipo}
\end{minipage}
\hfill
\begin{minipage}{0.48\textwidth}
\centering
\includegraphics[width=\textwidth]{5.4-historico-anticipos.png}
\caption{Histórico de anticipos}
\label{fig:historico-anticipos}
\end{minipage}
\end{figure}

\section{Finiquitos}

El finiquito es la liquidación final de las prestaciones sociales al momento de la desvinculación del servicio militar.

\subsection*{Triggers de Despido / Reserva Activa (Pase a Retiro)}

\begin{tcolorbox}[
    colback=sandraSoft,
    colframe=sandraDark,
    width=\textwidth,
    title=\textbf{Liquidación de Finiquito (Art. 57)}
]
\small
Toda fracción superior a \textbf{6 meses} equivale a \textbf{1 Año entero} frente a los dictados legales para la liquidación. Al generar el Finiquito total, el motor debe efectuar el cálculo en base a los 30 días continuos mensuales.

\begin{lstlisting}[language=Java, numbers=none]
// Suponiendo transformacion de meses extra a anosTotalesServicio
let montoPrestacionSocial = ultimoSueldoIntegralRetirado * anosTotalesServicio * 30;
\end{lstlisting}
\end{tcolorbox}

\subsection*{Tipos de Finiquito}

\begin{longtable}{p{3cm}p{9cm}}
\toprule
\textbf{Tipo} & \textbf{Descripción} \\
\midrule
\endhead
\bottomrule
\caption{Tipos de finiquito}
\label{tab:tipos-finiquito}
\endlastfoot

Retiro Voluntario & Desvinculación por decisión propia del afiliado \\
Retiro Obligatorio & Desvinculación por alcanzar límite de edad o servicio \\
Discapacidad & Jubilación por incapacidad física o mental \\
Fallecimiento & Liquidación a beneficiarios designados \\
Transferencia & Cambio de componente o situación militar \\
\end{longtable}

\subsection*{Componentes del Finiquito}

\begin{itemize}[leftmargin=*]
    \item \textbf{Aporte Patrimonial}: Acumulación de aportes del afiliado
    \item \textbf{Contribuciones Patronales}: Aportes del IPSFANB (cuota 2x)
    \item \textbf{Intereses Acumulados}: Rendimientos financieros
    \item \textbf{Ajustes por Inflación}: Compensación por variación monetaria
    \item \textbf{Deducciones}: Anticipos no liquidados, deudas, etc.
\end{itemize}

\section{Gestión Judicial (Art. 60 y 61)}

Este submódulo administra las medidas cautelares y embargos judicialmente ordenados.

\subsection*{Algoritmo Condicional de Indemnización por Muerte Súbita}

La diferencia radical algorítmica antes del pago obedece explícitamente a los estatutos 60 y 61:

\begin{lstlisting}[language=Java, caption=Indemnización de Muerte (Art. 60 y 61)]
let liquidacionTotal = 0;

switch (condicionFallecimientoStatus) {
    case 'ACTO_DE_SERVICIO':
        // Artículo 60: Máximo entre 36 sueldos militares o 36 salarios mínimos
        liquidacionTotal = Math.max(
            (ultimoSueldoMilitar * 36), 
            (salarioMinimoNacional * 36)
        );
        break;

    case 'FUERA_DE_SERVICIO':
        // Artículo 61: Máximo entre 24 sueldos militares o 24 salarios mínimos
        liquidacionTotal = Math.max(
            (ultimoSueldoMilitar * 24), 
            (salarioMinimoNacional * 24)
        );
        break;
}
\end{lstlisting}

\subsection*{Tipos de Medidas}

\begin{longtable}{p{3cm}p{9cm}}
\toprule
\textbf{Tipo} & \textbf{Descripción} \\
\midrule
\endhead
\bottomrule
\caption{Tipos de medidas judiciales}
\label{tab:medidas-judiciales}
\endlastfoot

Embargo & Prohibición de disposición de fondos por orden judicial \\
Retención & Deducción obligatoria de un porcentaje para terceros \\
Manutención & Descuento para pago de pensiones alimentarias \\
Decomiso & Incautación total o parcial por sentencia firme \\
\end{longtable}

\section{Movimientos del Fideicomiso}

Este submódulo muestra el historial completo de operaciones del afiliado:

\begin{figure}[H]
\begin{tcolorbox}[
    enhanced,
    colback=white,
    colframe=sandraDark,
    width=0.85\textwidth,
    halign=center
]
\begin{tabular}{cll}
\toprule
\textbf{Fecha} & \textbf{Concepto} & \textbf{Monto} \\
\midrule
15/01/2026 & Aporte Mensual & +\$\$500.000,00 \\
28/01/2026 & Interés Fideicomiso & +\$\$25.000,00 \\
05/02/2026 & Anticipo Educativo & -\$\$2.000.000,00 \\
15/02/2026 & Aporte Mensual & +\$\$500.000,00 \\
\midrule
\textbf{Saldo Actual} & & \textbf{\$\$8.525.000,00} \\
\bottomrule
\end{tabular}
\end{tcolorbox}
\caption{Ejemplo de movimiento de fideicomiso}
\label{fig:movimientos-ejemplo}
\end{figure}

\section{Órdenes de Pago}

Controla los disbursos realizados al afiliado:

\begin{itemize}[leftmargin=*]
    \item \textbf{Pendiente}: Orden creada, en espera de procesamiento
    \item \textbf{Ejecutada}: Pago realizado exitosamente
    \item \textbf{Fallida}: Error en el procesamiento del pago
    \item \textbf{Cancelada}: Orden anulada antes del procesamiento
\end{itemize}

\section{Procesamiento en Lote}

Permite la ejecución masiva de operaciones:

\begin{itemize}[leftmargin=*]
    \item \textbf{Generación de nóminas}: Creación batch de pagos
    \item \textbf{Liquidación de anticipos}: Procesamiento múltiple
    \item \textbf{Aplicación de intereses}: Cálculo y acreditación batch
    \item \textbf{Actualización de estados}: Cambio masivo de status
\end{itemize}

Este componente es de extrema delicadeza. Aquí el servidor itera sobre miles de Fideicomitentes para calcular (Módulo Cálculo) y abonar (Módulo Aportes).

\section{Reportes}

Generación de informes para consulta y archivo:

\begin{itemize}[leftmargin=*]
    \item \textbf{Estado de Cuenta}: Resumen detallado de operaciones
    \item \textbf{Certificado de Aportes}: Constancia de contribuciones
    \item \textbf{Historial de Anticipos}: Registro de solicitudes pasadas
    \item \textbf{Proyección de Jubilación}: Estimación de beneficios futuros
\end{itemize}

\begin{tcolorbox}[
    enhanced,
    colback=sandraSoft,
    colframe=sandraAccent,
    width=0.9\textwidth,
    halign=center
]
\small
Todos los reportes pueden ser exportados en formato PDF para impresión o archivo digital. También está disponible la exportación en formato Excel para análisis de datos adicionales.
\end{tcolorbox}

\newpage
